% Options for packages loaded elsewhere
\PassOptionsToPackage{unicode}{hyperref}
\PassOptionsToPackage{hyphens}{url}
\PassOptionsToPackage{dvipsnames,svgnames,x11names}{xcolor}
\documentclass[
  11pt,
  a4paper,
]{article}
\usepackage{xcolor}
\usepackage[margin=24mm]{geometry}
\usepackage{amsmath,amssymb}
\setcounter{secnumdepth}{5}
\usepackage{iftex}
\ifPDFTeX
  \usepackage[T1]{fontenc}
  \usepackage[utf8]{inputenc}
  \usepackage{textcomp} % provide euro and other symbols
\else % if luatex or xetex
  \usepackage{unicode-math} % this also loads fontspec
  \defaultfontfeatures{Scale=MatchLowercase}
  \defaultfontfeatures[\rmfamily]{Ligatures=TeX,Scale=1}
\fi
\usepackage{lmodern}
\ifPDFTeX\else
  % xetex/luatex font selection
\fi
% Use upquote if available, for straight quotes in verbatim environments
\IfFileExists{upquote.sty}{\usepackage{upquote}}{}
\IfFileExists{microtype.sty}{% use microtype if available
  \usepackage[]{microtype}
  \UseMicrotypeSet[protrusion]{basicmath} % disable protrusion for tt fonts
}{}
\makeatletter
\@ifundefined{KOMAClassName}{% if non-KOMA class
  \IfFileExists{parskip.sty}{%
    \usepackage{parskip}
  }{% else
    \setlength{\parindent}{0pt}
    \setlength{\parskip}{6pt plus 2pt minus 1pt}}
}{% if KOMA class
  \KOMAoptions{parskip=half}}
\makeatother
\usepackage{longtable,booktabs,array}
\newcounter{none} % for unnumbered tables
\usepackage{calc} % for calculating minipage widths
% Correct order of tables after \paragraph or \subparagraph
\usepackage{etoolbox}
\makeatletter
\patchcmd\longtable{\par}{\if@noskipsec\mbox{}\fi\par}{}{}
\makeatother
% Allow footnotes in longtable head/foot
\IfFileExists{footnotehyper.sty}{\usepackage{footnotehyper}}{\usepackage{footnote}}
\makesavenoteenv{longtable}
\setlength{\emergencystretch}{3em} % prevent overfull lines
\providecommand{\tightlist}{%
  \setlength{\itemsep}{0pt}\setlength{\parskip}{0pt}}
% Shared style; used by generated, standalone LaTeX sources.
\usepackage{xcolor}
\usepackage{xurl}
\usepackage{fvextra}
\usepackage{etoolbox}
\usepackage{fancyhdr}
\usepackage{titlesec}
\definecolor{researchblue}{HTML}{173B58}
\AtBeginDocument{\hypersetup{linkcolor=researchblue,urlcolor=researchblue}}
\pagestyle{fancy}
\fancyhf{}
\fancyhead[L]{\small\sffamily SVG PATCH LAB}
\fancyhead[R]{\small\sffamily CVPR 2027 research package}
\fancyfoot[L]{\small Working documents; results and proposals distinguished}
\fancyfoot[R]{\thepage}
\setlength{\headheight}{15pt}
\setlength{\emergencystretch}{3em}
\titleformat{\section}{\Large\bfseries\color{researchblue}}{\thesection}{0.7em}{}
\titleformat{\subsection}{\large\bfseries\color{researchblue}}{\thesubsection}{0.7em}{}
\DefineVerbatimEnvironment{verbatim}{Verbatim}{breaklines=true,breakanywhere=true,fontsize=\small}
\AtBeginEnvironment{longtable}{\small\setlength{\tabcolsep}{4pt}}
\let\originaltableofcontents\tableofcontents
\renewcommand{\tableofcontents}{\originaltableofcontents\clearpage}
\usepackage{bookmark}
\IfFileExists{xurl.sty}{\usepackage{xurl}}{} % add URL line breaks if available
\urlstyle{same}
\hypersetup{
  pdftitle={Research Proposal},
  pdfauthor={SVG Patch Lab - Research Working Package},
  colorlinks=true,
  linkcolor={Maroon},
  filecolor={Maroon},
  citecolor={Blue},
  urlcolor={blue},
  pdfcreator={LaTeX via pandoc}}

\title{Research Proposal}
\author{SVG Patch Lab - Research Working Package}
\date{14 September 2026 \textbar{} CVPR 2027 target}

\begin{document}
\maketitle

{
\setcounter{tocdepth}{2}
\tableofcontents
}
\section{Proposed title and status}\label{proposed-title-and-status}

\textbf{Physics-consistent editing of engineering vector drawings}

Target: CVPR 2027. This is a proposal backed by limited diagnostic
experiments, not an accepted contribution or completed submission.
Author list and affiliations are pending. The repository branch packages
the evidence, implementation, alternatives and remaining experiments
together.

\section{Research question}\label{research-question}

Can an editable engineering drawing preserve the meaning of its
numerical claims through natural-language presentation and physical
edits, while remaining useful for drafting?

A curve marked 40 degrees Celsius is a numerical assertion about every
point on that curve. A stress legend asserts a quantity, unit and scale.
Changing a load invalidates a previous solution even when its plot still
looks plausible. Changing an ancestor SVG transform can invalidate
geometry without changing the path string. The output to evaluate is the
final drawing, together with the physical problem and numerical evidence
supporting it.

\section{What prior work already
establishes}\label{what-prior-work-already-establishes}

MechAgents, MCP-SIM and OASiS connect language models to numerical
simulation. FEABench and PDEAgent-Bench evaluate simulation and solver
generation. FEM-Bench evaluates FEM functions and tests. ALL-FEM and
ModSolAgent already fine-tune models for engineering solver workflows.
VFEAgent addresses image/text-to-FEA. Penrose separates mathematical
semantics from visual style; scientific SVG generation and editing have
their own benchmarks and methods.

Therefore, LLM plus FEM, verified-code fine-tuning, multi-agent
orchestration, and multimodal input are not sufficient novelty claims.
The proposed contribution concerns consistency between the solved
physical problem and the final editable artifact, particularly across a
sequence of edits. The companion literature survey gives source links,
reuse choices and known release limitations. It does not prove that the
proposed combination is absent from all prior work.

\section{Proposed representation}\label{proposed-representation}

Each artifact has four associated records:

\begin{enumerate}
\def\labelenumi{\arabic{enumi}.}
\tightlist
\item
  \textbf{Physical specification:} domain, dimensional assumptions, PDE,
  material law, loads, supports, boundary conditions, units and
  requested quantities. Unknown values remain explicit.
\item
  \textbf{Numerical solution:} solver and version, discretization, mesh,
  tolerances, fields, reactions/fluxes, convergence checks and
  provenance.
\item
  \textbf{Drawing claims:} SVG element IDs, physical-coordinate mapping,
  field quantity, contour level, associated text, unit and legend
  binding.
\item
  \textbf{Edit record:} instruction, parsed operation, affected physical
  inputs or presentation properties, numerical invalidation, solver
  execution and verifier outcome.
\end{enumerate}

Content hashes record identity; they do not prove that a solution is
correct. A reference solution must pass its own physical and numerical
checks. The representation must distinguish a computed field, a sampled
geometric check, a schematic illustration and an unsupported claim.

\section{Proposed system}\label{proposed-system}

The planner resolves the engineering specification and requests missing
information only when necessary for a well-posed problem. A solver
executes the specified model. A deterministic exporter creates the
numerical field geometry. The editor proposes changes to geometry,
annotations or presentation through a structured interface. The verifier
checks the final artifact and either accepts it within a stated scope,
reports a limitation, or requests a repair/re-solve.

Presentation edits preserve the physical solution but may change layout,
line style and annotation position. Physical edits create a new
specification and trigger a new solve. Unsupported semantic changes
cannot retain the old verification status. Requests that contradict the
specification must be rejected or clearly marked as illustrative.

For the main study, the system must inspect visible labels, units,
legends and occlusion in addition to stored metadata. Merely attaching a
solution ID to an SVG is insufficient. Both useful layout flexibility
and numerical/semantic consistency must be measured.

\section{Hypotheses and
falsification}\label{hypotheses-and-falsification}

{\def\LTcaptype{none} % do not increment counter
\begin{longtable}[]{@{}
  >{\raggedright\arraybackslash}p{(\linewidth - 6\tabcolsep) * \real{0.2500}}
  >{\raggedright\arraybackslash}p{(\linewidth - 6\tabcolsep) * \real{0.2500}}
  >{\raggedright\arraybackslash}p{(\linewidth - 6\tabcolsep) * \real{0.2500}}
  >{\raggedright\arraybackslash}p{(\linewidth - 6\tabcolsep) * \real{0.2500}}@{}}
\toprule\noalign{}
\begin{minipage}[b]{\linewidth}\raggedright
ID
\end{minipage} & \begin{minipage}[b]{\linewidth}\raggedright
Hypothesis
\end{minipage} & \begin{minipage}[b]{\linewidth}\raggedright
Evidence required
\end{minipage} & \begin{minipage}[b]{\linewidth}\raggedright
Result that would weaken it
\end{minipage} \\
\midrule\noalign{}
\endhead
\bottomrule\noalign{}
\endlastfoot
H1 & Joint geometric and semantic checking reduces false acceptance
after edits. & Fewer accepted wrong-field, wrong-unit and stale-solution
outputs at comparable edit success. & Geometry-only or deterministic
controls perform equally well. \\
H2 & A structured physical-edit path prevents reuse of stale fields. &
Actual solver reruns and correct updated plots after load, material,
boundary and geometry changes. & Provenance IDs change but plotted
numbers remain stale. \\
H3 & Flexible protected editing is more useful than freezing the
drawing. & Blinded drafting and instruction-fulfillment ratings improve
without weakening numerical checks. & A simple plotting template matches
quality and success. \\
H4 & Learning helps on context-dependent, visually grounded edits. & SFT
improves over the same base model, executor and information on held-out
cases/templates. & A rule parser or prompting closes the gap. \\
H5 & The method generalizes beyond a single PDE or SVG template. &
Independent geometry-family, instruction-template and sequential-edit
holdouts. & Gains disappear when templates or layouts change. \\
\end{longtable}
}

The study should report negative results. A strong deterministic
baseline is useful evidence about the necessity of learned components,
not a baseline to weaken artificially.

\section{Evidence already available}\label{evidence-already-available}

The repaired contour evaluator samples curves and composes ancestor
transforms, retains unsupported-reference samples, and checks required
levels. Thirty regression tests pass. It remains a sampled geometric
audit; arbitrary labels, units, legends and occlusion are not fully
covered.

On 160 manufactured controls, the geometric checker accepts 20 of 120
corrupted examples, all wrong-label cases. A rigid freeze baseline
accepts none. This identifies missing semantic checks, but does not
establish superiority over the rigid baseline.

Re-auditing six historical drawings gives three strict geometric passes.
The other outcomes include incompatible-corner reference ambiguity,
masked boundary coverage and an empty response. They are not three
demonstrated incorrect physical solutions.

Four fresh Astra API calls revisit the suspected failures. Both Robin
heat drawings pass direct FEM contour sampling, with maximum errors
0.111 and 0.374 degrees Celsius. The old closed-loop claim about the
capacitor is contradicted by its saved paths. The bracket discloses its
illustrative stress field. These corrections rule out an unqualified
claim that these cases demonstrate general FEM incapability.

Three A100 LoRA runs completed on the compact-DOM action pilot. Every
final adapter and the deterministic rule parser score 60/60 on both
evaluation splits. This is training infrastructure validation, not
evidence of a new engineering or visual capability.

A working extension imports the unchanged FEM-Bench axial-bar solver and
performs 24 actual FEM solves, including load edits. It passes the two
upstream reference tests, analytic displacement and force-balance
checks. It is an elementary integration anchor, not the main benchmark.

\section{Main benchmark proposal}\label{main-benchmark-proposal}

Begin with well-posed Robin conduction, finite-domain linear elasticity
and electrostatics. Each family needs analytic anchors, converged
numerical references and seeded geometry/material variants. Specify
plane stress versus plane strain or 3-D assumptions, traction
distributions, supports, fillets and far-boundary treatment. Do not
grade an infinite-plate formula as exact for a finite square plate.

Keep all variants of a physical case in one split. Separate unseen
geometry from unseen language templates. Include honest counterfactuals
and misleading requests, as well as ordinary restyling and layout
changes. Multi-step sequences should alternate physical changes and
presentation changes so that stale claims can be exposed.

Compare direct Astra editing, unrestricted solver-informed editing, a
published solver-agent workflow, protected geometry, joint protection
and deterministic solver plotting. Use equal numerical information where
the comparison requires it. Human evaluation should assess readability,
instruction fulfillment and visible claim correctness on anonymized,
randomized outputs.

\section{Training proposal}\label{training-proposal}

Do not scale the current perfect-template pilot as the main experiment.
First establish a difficult task and an improvement opportunity beyond
deterministic and prompt-only controls. Construct training examples from
verified solver records and successful edits, with rejected edits and
corrective feedback represented separately.

Use the current small open-weight model as an initial cost-bounded
baseline. Add image/complete-SVG context only with a model and
representation that actually support those inputs. Freeze the evaluation
protocol before hyperparameter selection. Use validation for selection
and three seeds for a final justified configuration. Test-time cases,
failed-example inspection and human instructions used in evaluation must
not silently enter training.

\section{Alternative proposals}\label{alternative-proposals}

\textbf{A. Benchmark and audit contribution.} Build a carefully
calibrated benchmark of numerical/semantic failures in engineering
drawings. This can remain valuable if no new learned method is needed.
Its strength would be reference quality, realistic editing tasks and
reproducible assessment, not a large synthetic sample count alone.

\textbf{B. Solver-linked graphical representation.} Focus on the
representation and verifier, comparing useful edits to rigid freezing
and deterministic plotting. This is the preferred initial direction
because it connects the numerical and graphical parts of the existing
implementation.

\textbf{C. Learned solver-code assistant.} Extend ALL-FEM or FEM-Bench
with training on engineering programs. This overlaps strongly with
existing work and has a weaker CVPR fit unless a substantive visual
component is established. Do not present it as a novel fallback merely
because GPU training is available.

\textbf{D. Learned numerical surrogate.} Train a PINN or field surrogate
for a specific parametric family and compare its total cost and accuracy
with FEM. This is a separate scientific question and should not be added
solely to include a neural solver.

\section{Risks and decision gates}\label{risks-and-decision-gates}

The principal risks are reference error masquerading as model error,
under-specified engineering tasks, trivial templated edits, visual
quality lost through over-restriction, and overlap with prior systems.
Address these through independent checks, explicit task specifications,
strong deterministic controls, blinded evaluation and a full baseline
audit.

Gate 1: references and elementary integration pass before main model
comparisons. Gate 2: the benchmark exposes a useful gap before large
training. Gate 3: a visual/graphics contribution survives the strongest
controls before claiming CVPR readiness. If these gates fail, narrow the
claim or change venue rather than inflate the result.

\section{Delivery and schedule}\label{delivery-and-schedule}

The official CVPR 2027 dates currently list registration on 10 November
2026, main paper on 16 November and supplementary material on 23
November, all Anywhere on Earth. Recheck the current call and author
instructions before submission:
\url{https://cvpr.thecvf.com/Conferences/2027/Dates}.

September: lock references, claim schema, realistic task protocol and
baseline environments. October: main comparisons, justified training,
sequential edits, ablations and human assessment. Early November:
complete results, anonymous paper, figures, limitations and artifact
instructions. Submission requires authorship information, author review
and final approval; none is implied by this package.

Deliverables include source code, exact experiment configurations,
input/reference/output records, per-case metrics, compute accounting,
licenses, paper sources and a reproducibility guide. The companion
experiment protocol translates these goals into explicit jobs and
acceptance criteria.

\end{document}
